% !TeX root = ../main.tex
\chapter{Conclusion and Future Work}

\section{Summary of Contributions}

This thesis set out to bridge the gap between the theoretical capacity of Mixture-of-Experts (MoE) architectures and the strict computational constraints of on-device chess evaluation functions. By moving away from traditional, handcrafted bucketing based on piece count, we explored data-driven routing mechanisms capable of unlocking deeper expert specialization. The primary contribution is the formalization of the "routing gap" in NNUE architectures: the fundamental dichotomy between the feature space that is easily dispatchable at inference time (L1 activations) and the feature space that contains the true specialization signal (sample gradients).

Through extensive empirical evaluation, we established that while a strong reference base model leaves little room for single-stage MoE improvements, the underlying experts possess significant latent capacity. Our oracle upper bounds proved that if routed perfectly, the experts could achieve substantial cross-entropy reductions. Furthermore, we demonstrated that when applied to a slightly weaker base architecture, L1-clustered hard MoE variants could indeed outperform the baseline, provided the experts were given sufficient training data. This confirmed that the failure of MoE to beat the strong base was not a lack of capacity, but a failure of the inference-time routing signal to recover the optimal expert assignment.

A critical secondary contribution was the identification and resolution of the optimization instabilities inherent in fine-tuning multiple expert heads. By transitioning to the Rprop optimizer and implementing a dynamic learning rate adjustment that steps back and slows down when the loss increases, we achieved a highly stable and reliable fine-tuning phase. This optimization breakthrough was instrumental in allowing the experts to properly specialize without collapsing or overfitting, unlocking the first reproducible routing gains in our experimental battery.

To resolve the routing gap without sacrificing inference efficiency, this thesis proposed KITSUNE, a two-stage routing framework. By decoupling the geometric partitioning of the L1 space into micro-buckets from the offline, gradient-based stitching of those buckets into macro-experts, KITSUNE lets a linear dispatcher express non-convex gradient regions. On the weaker \texttt{W64 H128} base, the routed mixture at $K=16$ and factor $2$ ($B'=32$ micro-buckets) reaches CE 0.65043, a reduction of 0.0033 from the base CE of 0.6537 and about 0.0005 above the best single-stage L1 mixture. The best expert on its own bucket reaches CE 0.60571 at that setting (a 7.34\% reduction against the base CE) and CE 0.60234 at factor $4$. Gradient-informed stitching therefore produces specialized heads while the online router stays a linear map plus a lookup. The deployed mixture does not yet surpass the best single-stage L1 run, and factor $8$ at $K=16$ worsens both the mixture and the best expert.

\section{Future Work}

The $K \in \{4, 8, 16\}$ sweep places the useful micro-bucket ratios at $B'/K \in \{2, 4\}$; factor $8$ at $K=16$ already degrades both the routed mixture and the best expert. The immediate next step is a load-balanced stitching rule, so that no macro-expert absorbs most of the data, together with a direct comparison of the same KITSUNE configuration against the strong \texttt{W128 H256} reference. That comparison is still open: every KITSUNE number in Chapter 5 was measured on the weaker base.

While cross-entropy and MAE are excellent proxies for evaluation quality, they do not perfectly correlate with playing strength in a full search engine. Future work must involve integrating the KITSUNE dispatcher and the specialized expert heads directly into the C++ search engine. This will require optimizing the sparse matrix operations for the specific instruction sets of modern CPUs and measuring the actual nodes-per-second (NPS) overhead. Ultimately, the true value of the architecture will be validated through Elo rating tournaments against both the baseline NNUE and existing handcrafted bucket implementations.

The same split is a candidate for other shared-trunk models, including multi-task networks and recommendation systems. Carrying it over means re-fitting the micro-buckets and the gradient stitch on that domain's trunk and loss. The chess sweep does not establish that the split transfers.

\section{Closing Remarks}
The pursuit of highly specialized evaluation functions for constrained environments often forces a compromise between model capacity and inference speed. This thesis demonstrates that this compromise is not absolute. By rigorously analyzing the geometry of the routing space and introducing a two-stage, offline-computed dispatching mechanism, we can extract the deep specialization benefits of Mixture-of-Experts without incurring the computational penalties of dynamic, non-linear routing at runtime. Data-driven, gradient-informed partitions can replace a handcrafted piece-count heuristic, and a linear dispatcher plus an offline stitch is enough to keep that partition deployable on constrained hardware.